\documentclass[11pt]{article}
\usepackage[margin=1in]{geometry}
\usepackage{amsmath,amssymb,amsthm}
\usepackage{booktabs}
\usepackage{hyperref}
\usepackage{enumitem}

\newtheorem{theorem}{Theorem}
\newtheorem{corollary}[theorem]{Corollary}
\newtheorem{lemma}[theorem]{Lemma}
\newtheorem{proposition}[theorem]{Proposition}
\newtheorem{definition}[theorem]{Definition}
\newtheorem{remark}[theorem]{Remark}

\title{Frozen Core Isolation and Quantum-Resistant\\Cryptographic Commitments from Planted $k$-SAT}
\author{John Rhodes}
\date{}

\begin{document}
\maketitle

\begin{abstract}
We prove that for planted $k$-SAT instances with $k \geq 7$ at clause density $\alpha/\alpha_s \geq 0.21$, a positive fraction of variables are frozen directly in the planted model---without requiring transfer from the random model via quiet planting. The expected number of ``support clauses'' per variable (clauses in which that variable is the unique satisfying literal) exceeds~1 at remarkably low density: $\alpha/\alpha_s \approx 0.20$ for $k = 7$, compared to the random-model freezing threshold at $\alpha_f/\alpha_s \approx 0.90$. We prove that the resulting frozen-core structure implies topological disconnection of the solution subgraph across cluster boundaries, with a cycle-robustness argument showing that short cycles in the factor graph cannot quench the supercritical repair cascade. As an immediate corollary, the Hilbert space spanned by satisfying assignments decomposes into orthogonal sectors preserved by any unitary generated by the adjacency matrix---blocking quantum walks, QAOA at all depths, and quantum annealing. We construct a post-quantum commitment scheme whose binding property reduces to the hardness of solving planted $k$-SAT, provide formal proofs of completeness, soundness, and zero-knowledge, and derive a digital signature scheme with existential unforgeability via the Fiat--Shamir transform.  We present a six-vector quantum attack analysis with proved barriers against five algorithmic families. We give concrete parameter recommendations at NIST security levels 1, 3, and~5, and position the scheme within the landscape of SAT-based and CSP-based cryptographic constructions. We prove that the Grover query complexity for breaking the binding property is $\Omega(2^{fn/2})$; empirical cryptanalysis of Glucose and MiniSat CDCL solvers on our exact distribution yields a classical attack cost of $2^{0.234n}$ operations, enabling concrete parameter selection at NIST security levels 1, 3, and~5. Empirical validation across 100 random seeds at $n = 16$ confirms complete cluster isolation at every instance tested.
\end{abstract}

%% ===================================================================
\section{Introduction}
%% ===================================================================

Post-quantum cryptography rests on a narrow foundation.  The three hardness families underlying deployed schemes---lattices, hash functions, and error-correcting codes---represent a concentration of risk: a breakthrough against any one family would compromise a large fraction of the post-quantum infrastructure.  Diversifying the set of available hardness sources is a security imperative~\cite{BL17}. This concern is sharpened by recent work proposing additional hardness families---quantum stabilizer decoding~\cite{LPQR26} and, in the present paper, frozen-core constraint satisfaction---suggesting the viable assumption space is broader than the current standards reflect.

We introduce a fourth hardness family based on the cluster structure of random constraint satisfaction problems.  Our main contributions are:

\begin{theorem}[Frozen Core Isolation]\label{thm:fci}
For planted $k$-SAT with $k \geq 7$ and clause density $\alpha$ satisfying $e_{\mathrm{sup}}(k,\alpha) > 1$ (equivalently, $\alpha/\alpha_s > 1/\bigl(k \cdot 2^k \ln 2 \cdot (2^k - 1)^{-1}\bigr) \approx 0.20$ for $k = 7$), the following hold with high probability:
\begin{enumerate}[label=(\roman*)]
\item Each variable has expected support count $e_{\mathrm{sup}} = k\alpha/(2^k - 1) > 1$, where a support clause is one in which the variable is the unique satisfying literal under the planted solution~$y^*$.
\item A fraction $f = 1 - e^{-e_{\mathrm{sup}}(1-q)} > 0$ of variables are frozen, where $q < 1$ is the extinction probability of a Galton--Watson process with offspring mean $e_{\mathrm{sup}}$.
\item The connected component $\mathrm{Comp}(y^*)$ of $y^*$ in the Hamming-1 solution subgraph $G_{\mathrm{SAT}}$ is contained in $B(y^*,(1-f+\varepsilon)n)$ for any $\varepsilon > 0$.
\item For two planted solutions $y_1, y_2$ with $\mathrm{Ham}(y_1,y_2) > 2(1-f)n$, the connected components $\mathrm{Comp}(y_1)$ and $\mathrm{Comp}(y_2)$ are disjoint.
\end{enumerate}
\end{theorem}

\begin{corollary}[Quantum Isolation]\label{cor:qi}
Under the hypotheses of Theorem~\ref{thm:fci}, the Hilbert space $\mathcal{H} = \mathrm{span}\{|s\rangle : s \text{ satisfies } \Phi\}$ decomposes as $\mathcal{H} = \bigoplus_i \mathcal{H}_i$, where each $\mathcal{H}_i$ corresponds to a connected component of $G_{\mathrm{SAT}}$.  Any unitary of the form $g(A)$, where $A$ is the adjacency matrix of $G_{\mathrm{SAT}}$ and $g$ is any analytic function (including $e^{-iAt}$), preserves this decomposition.
\end{corollary}

\begin{theorem}[Commitment and Signature Schemes]\label{thm:commit}
There exists a post-quantum commitment scheme with perfect completeness, soundness error $2^{-\lambda}$ (for security parameter~$\lambda$), and computational zero-knowledge, whose binding property reduces to the hardness of solving planted $k$-SAT.  The underlying sigma protocol yields, via the Fiat--Shamir transform, a digital signature scheme with existential unforgeability under chosen-message attack (EUF-CMA) in the random oracle model.
\end{theorem}

The paper does not claim quantum circuit lower bounds for planted $k$-SAT.  No existing post-quantum scheme has such a proof.  What we provide is, in certain respects, more explicit quantum resistance evidence than is available for existing schemes: proved barriers against five of six major quantum algorithmic families (Table~\ref{tab:quantum}), with only generic brute-force search (Grover) remaining as an unstructured threat mitigated by standard parameter doubling.


%% ===================================================================
\subsection{Related work and positioning}\label{sec:related}
%% ===================================================================

\paragraph{Goldreich's one-way function.}
Goldreich~\cite{Gol00} proposed planted constraint satisfaction problems as candidates for one-way functions.  His construction evaluates a fixed $d$-ary predicate over an expander graph; security evidence comes from exponential lower bounds against myopic backtracking algorithms~\cite{CEMT09} and statistical algorithm lower bounds~\cite{FPV18}.  Our construction differs in three ways: (a)~we use standard planted $k$-SAT rather than a custom predicate on an expander, (b)~our hardness mechanism is topological---frozen-core cluster isolation---rather than algorithmic-class-specific, and (c)~we provide explicit quantum resistance analysis (Table~\ref{tab:quantum}) with no counterpart in the Goldreich OWF literature.

\paragraph{Feldman--Perkins--Vempala.}
Feldman, Perkins, and Vempala~\cite{FPV18} proved tight $\Omega(n^{r/2})$ lower bounds for statistical algorithms on planted $k$-SAT, matching known upper bounds.  Their results provide evidence for the hardness of \emph{finding} the planted solution.  Our Theorem~\ref{thm:fci} addresses a complementary question: even if an adversary finds \emph{some} satisfying assignment, the frozen-core structure guarantees it lies in the same cluster as the planted solution, making the binding property of our commitment scheme structurally enforced.  The FPV statistical algorithm bounds strengthen our hardness assumption by ruling out a broad class of classical attacks.

\paragraph{SAT-based cryptography.}
Schmittner~\cite{Sch15} proposed a public-key encryption scheme using a random SAT formula as public key and a satisfying assignment as private key.  That construction lacks any structural hardness analysis---no planted model, no freezing threshold computation, no cluster isolation, and no quantum security argument.  More recently, Cho et~al.~\cite{Cho25} proposed ``Equilibrium SAT Based PQC'' using a counting/multiset mechanism over SAT instances.  Neither work exploits the phase-transition structure of random CSPs; both treat SAT as a generic NP-hard problem.  Our approach is fundamentally different: hardness derives not from the generic NP-completeness of SAT but from the specific geometric structure (frozen clusters, topological disconnection) that emerges at controlled clause densities in the planted model.

\paragraph{Frozen variables and hardness.}
The connection between frozen variables and computational hardness was conjectured by Zdeborov\'{a} and Krzakala~\cite{ZK07} and studied rigorously for random CSPs by Achlioptas and Ricci-Tersenghi~\cite{ART09}, Molloy~\cite{Mol12}, and Coja-Oghlan and Zdeborov\'{a}~\cite{CZ12}.  All prior analysis concerns the \emph{random} model.  Our Lemma~\ref{lem:support} provides, to our knowledge, the first computation of the freezing threshold in the \emph{planted} model, revealing that planted instances freeze at roughly $1/4.4$ the density required in the random model (for $k = 7$).  Munoz-Bauza et~al.~\cite{MMRE23} explicitly identified the formal connection between frozen variables and algorithmic hardness as an open problem in the context of planted Ising instances; our Theorem~\ref{thm:fci} resolves this for planted $k$-SAT.

\paragraph{New PQC hardness families.}
Lu, Poremba, Quek, and Ramkumar~\cite{LPQR26} recently proposed quantum stabilizer decoding as a new post-quantum hardness family, arguing it is plausibly incomparable to LPN.  Our work is parallel in motivation---diversifying the PQC hardness landscape beyond lattices, hashes, and codes---but based on entirely different mathematics.  Together these efforts suggest that the space of viable PQC assumptions is broader than the current NIST portfolio reflects.


%% ===================================================================
\section{Preliminaries}
%% ===================================================================

\subsection{Planted $k$-SAT}
Fix integers $n$ (variables) and $k$ (clause width).  A planted $k$-SAT instance is generated as follows:
\begin{enumerate}
\item Choose a planted solution $y^* \in \{0,1\}^n$ uniformly at random.
\item Generate $m = \lfloor \alpha n \rfloor$ clauses independently: each clause selects $k$ distinct variables uniformly and $k$ literal signs uniformly, conditioned on $y^*$ satisfying the clause.
\end{enumerate}
The satisfiability threshold is $\alpha_s \approx 2^k \ln 2$~\cite{DSS22}.

\subsection{Cluster structure and frozen variables}
For $\alpha$ above the clustering threshold $\alpha_d$, the set of satisfying assignments decomposes into exponentially many well-separated clusters~\cite{DSS22,ACO08}.  Above the freezing threshold $\alpha_f > \alpha_d$, each cluster contains a positive fraction of frozen variables~\cite{ART09,CZ12}.

\begin{definition}[Frozen variable]\label{def:frozen}
Variable $x_i$ is frozen in cluster $C$ if $x_i$ takes the same value in every solution reachable from $C$'s core solution by a path of single-flip satisfying assignments.  Equivalently, $x_i$ is constant on the connected component of $C$'s core in $G_{\mathrm{SAT}}$.
\end{definition}

\begin{definition}[Solution subgraph]\label{def:gsat}
$G_{\mathrm{SAT}} = (S, E)$ has vertex set $S = \{s \in \{0,1\}^n : s \text{ satisfies } \Phi\}$ and edge set $E = \{\{s,t\} : \mathrm{Ham}(s,t) = 1\}$.
\end{definition}


%% ===================================================================
\section{Proof of Theorem~\ref{thm:fci}}
%% ===================================================================

The proof proceeds in five steps: (1)~compute the support count directly in the planted model, (2)~establish tree-likeness of the planted factor graph, (3)~prove that supercritical support implies freezing via Warning Propagation, (4)~derive Hamming ball containment, and (5)~conclude component disjointness.

\subsection{Step 1: Support count in the planted model}

This is the key technical contribution.  We analyze the planted clause distribution directly, without transferring results from the random model.

\begin{lemma}[Planted support count]\label{lem:support}
In a planted $k$-SAT instance at density $\alpha$, the expected number of support clauses for a variable $x_i$ (clauses in which $x_i$ is the unique satisfying literal under $y^*$) is
\[
e_{\mathrm{sup}} = \frac{k\alpha}{2^k - 1}.
\]
\end{lemma}

\begin{proof}
Consider a clause $c$ containing variable $x_i$.  The number of clauses containing $x_i$ is $\mathrm{Bin}(m, k/n)$, with expectation $k\alpha$.

For such a clause, the literal for $x_i$ has sign $+1$ or $-1$ with equal probability before conditioning on acceptance (i.e., $y^*$ satisfying the clause).  Let $p_{\mathrm{acc}}$ denote the acceptance probability.

If $x_i$ appears with sign matching $y^*_i$ (``positive''): the clause is always accepted, since $x_i$ provides one satisfying literal.  Probability of this sign: $1/2$.  Conditional acceptance: $1$.

If $x_i$ appears with sign opposite to $y^*_i$ (``negative''): the clause is accepted only if at least one of the other $k-1$ literals is satisfied by $y^*$.  Probability: $1/2$.  Conditional acceptance: $1 - 2^{-(k-1)}$.

Thus $p_{\mathrm{acc}} = \frac{1}{2} \cdot 1 + \frac{1}{2}(1 - 2^{-(k-1)}) = 1 - 2^{-k}$.

Given that $c$ is accepted and contains $x_i$, the probability that $x_i$ is the unique satisfying literal is:
\[
P(\text{unique}) = \frac{P(x_i \text{ positive} \mid \text{acc}) \cdot P(\text{all other } k{-}1 \text{ lits false} \mid x_i \text{ pos})}{1} = \frac{1/2}{1 - 2^{-k}} \cdot 2^{-(k-1)} = \frac{1}{2^k - 1}.
\]
Multiplying by the expected number of clauses containing $x_i$:
\[
e_{\mathrm{sup}} = k\alpha \cdot \frac{1}{2^k - 1}. \qedhere
\]
\end{proof}

\begin{remark}[Planted vs.\ random freezing threshold]\label{rem:threshold}
Setting $e_{\mathrm{sup}} = 1$ gives the planted freezing threshold $\alpha_f^{\mathrm{plant}} = (2^k - 1)/(k \cdot 2^k \ln 2)$.  For $k = 7$: $\alpha_f^{\mathrm{plant}}/\alpha_s \approx 0.204$.  The random-model freezing threshold is $\alpha_f^{\mathrm{rand}}/\alpha_s \approx 0.898$~\cite{ART09,CZ12}.  The planted model freezes at $4.4$ times lower density than the random model.  At our operating point $\alpha/\alpha_s = 0.78$: $e_{\mathrm{sup}} \approx 3.81$, deeply supercritical.
\end{remark}

\begin{table}[ht]
\centering
\caption{Planted freezing threshold $\alpha_f^{\mathrm{plant}}/\alpha_s$ vs.\ random-model estimate $\alpha_f^{\mathrm{rand}}/\alpha_s$.}\label{tab:freeze}
\begin{tabular}{@{}ccccc@{}}
\toprule
$k$ & $e_{\mathrm{sup}}$ formula & $\alpha_f^{\mathrm{plant}}/\alpha_s$ & $\alpha_f^{\mathrm{rand}}/\alpha_s$ & Ratio \\
\midrule
5  & $5\alpha/31$    & 0.280 & 0.719 & 0.39 \\
7  & $7\alpha/127$   & 0.205 & 0.898 & 0.23 \\
9  & $9\alpha/511$   & 0.160 & 0.967 & 0.17 \\
15 & $15\alpha/32767$ & 0.096 & 0.999 & 0.10 \\
\bottomrule
\end{tabular}
\end{table}


\subsection{Step 2: Tree-likeness of the planted factor graph}

\begin{lemma}[Planted tree-likeness]\label{lem:tree}
The factor graph of planted $k$-SAT at density $\alpha < \alpha_s$ has the same vertex-edge topology as a uniformly random $k$-uniform hypergraph with $m = \alpha n$ hyperedges on $n$ vertices.  In particular, the depth-$r$ neighborhood of any variable node is a tree with probability $1 - o(1)$ for $r = o\bigl(\sqrt{\log n / \log(k\alpha)}\bigr)$.
\end{lemma}

\begin{proof}
In the planted model, each clause selects $k$ variables uniformly at random and then selects literal signs conditioned on $y^*$ satisfying the clause.  The conditioning affects signs only, not which variables appear in which clause.  The hypergraph structure (clause-variable incidence) is therefore identically distributed to the standard random $k$-uniform hypergraph.  The result follows from the standard local tree-likeness theorem for random hypergraphs~\cite{DM10,MR95}.
\end{proof}

\begin{remark}
At the densities relevant to our construction ($k\alpha \approx 484$ for $k = 7$, $\alpha/\alpha_s = 0.78$), the tree-like depth is $r = O(1)$.  This suffices for Warning Propagation, which converges in $O(1)$ iterations when the clause-to-variable ratio exceeds the freezing threshold~\cite{CZ12}.
\end{remark}


\subsection{Step 3: Supercritical support implies freezing}

\begin{lemma}[Frozen--disconnected]\label{lem:frozen}
For planted $k$-SAT at density $\alpha$ with $e_{\mathrm{sup}} > 1$, let $F$ denote the set of variables receiving at least one hard warning under Warning Propagation on the factor graph.  Then w.h.p.:
\begin{enumerate}[label=(\alph*)]
\item $|F| \geq fn$ for $f = 1 - e^{-e_{\mathrm{sup}}(1-q)} > 0$, where $q < 1$ is the extinction probability of a Galton--Watson process with offspring distribution $\mathrm{Poi}(e_{\mathrm{sup}})$.
\item For each $x_i \in F$, the value $x_i = y^*_i$ is constant on $\mathrm{Comp}(y^*)$.
\end{enumerate}
\end{lemma}

\begin{proof}[Proof of (a)]
On the tree-like factor graph (Lemma~\ref{lem:tree}), Warning Propagation has a unique fixed point.  A variable $x_i$ receives a hard warning from clause $c$ if $x_i$ is the unique satisfying literal in $c$ under $y^*$.  By Lemma~\ref{lem:support}, the number of hard warnings per variable converges to $\mathrm{Poi}(e_{\mathrm{sup}})$ on the local tree.  The fraction of variables receiving $\geq 1$ hard warning is $1 - e^{-e_{\mathrm{sup}}} > 0$.

Not every hard warning produces a frozen variable; $x_i$ is frozen only if the repair cascade triggered by flipping $x_i$ diverges.  The cascade is a Galton--Watson process with offspring mean $e_{\mathrm{sup}}$: each flipped variable breaks its own support clauses, each of which requires repairing a different variable.  For $e_{\mathrm{sup}} > 1$ the process is supercritical, with survival probability $1 - q > 0$.  The fraction of variables for which at least one warning triggers a surviving cascade is $f = 1 - e^{-e_{\mathrm{sup}}(1-q)}$.

Concentration follows by McDiarmid's inequality: the frozen status of each $x_i$ depends on the $O(1)$ clauses in its depth-1 neighborhood, and changing a single clause affects at most $k$ variables' frozen status.  Setting $F = |\{i : x_i \text{ frozen}\}|$ as a function of $m$ independent clause choices, each with bounded difference $k$:
\[
\Pr\bigl[|F - fn| > n^{2/3}\bigr] \leq 2\exp\Bigl(-\frac{2n^{4/3}}{mk^2}\Bigr) = 2\exp\bigl(-\Omega(n^{1/3})\bigr) \to 0.
\]
Therefore $|F| = fn \pm o(n)$ w.h.p.
\end{proof}

\begin{proof}[Proof of (b)]
Suppose for contradiction that $s \in \mathrm{Comp}(y^*)$ satisfies $s_i \neq y^*_i$ for some $x_i \in F$.  Then there exists a path $y^* = s_0, s_1, \ldots, s_T = s$ in $G_{\mathrm{SAT}}$.  Let $t^*$ be the first index at which $x_i$ is flipped.  At step $t^*$, all support clauses for $x_i$ become violated.  Each must be repaired by flipping another variable in a subsequent step, but each repair flip breaks that variable's own support clauses.

\emph{First-passage property.}  Any repair sequence flips each variable at most once: a second flip of $x_j$ would re-break the clause whose violation caused the first flip of~$x_j$.  Therefore the cascade defines a tree~$T$ of distinct flipped variables, regardless of cycles in the factor graph.

\emph{Cycle encounters are automatic repairs.}  Suppose the cascade reaches a support clause~$c$ for some variable~$v$, and one of the $k-1$ other variables in~$c$, say~$x_j$, was already flipped earlier in the cascade.  Since $c$ is a support clause for~$v$, every literal in~$c$ other than~$v$'s is false under~$y^*$; in particular, $x_j$ appears in~$c$ with the sign opposite to~$y^*_j$.  Flipping~$x_j$ from~$y^*_j$ to~$\neg y^*_j$ therefore makes $x_j$'s literal in~$c$ true, satisfying~$c$ without any additional flip.

\emph{Stochastic domination.}  Each such automatic repair removes one branch from the cascade tree~$T$ (the sub-cascade that would have followed from flipping a fresh repair variable for~$c$).  The cascade on the full graph is therefore a strict pruning of the cascade on the local tree: every cycle encounter removes branches, never adds them.  The extinction probability satisfies $q_{\mathrm{graph}} \geq q_{\mathrm{tree}} = q$.

\emph{Cycle sparsity.}  Define the directed cascade graph~$G_C$: variable $x_a$ points to variable $x_b$ if there exists a support clause for~$x_a$ containing~$x_b$.  Each variable has expected out-degree $e_{\mathrm{sup}}(k-1)$ in~$G_C$.  The probability that a specific pair $(x_a, x_b)$ shares a cascade edge is $e_{\mathrm{sup}}(k-1)/n$.  A directed cycle of length~$\ell$ through~$\ell$ specified variables occurs with probability $(e_{\mathrm{sup}}(k-1)/n)^\ell$.  The expected number of directed $\ell$-cycles is $(e_{\mathrm{sup}}(k-1))^\ell / \ell$, and the expected number of such cycles \emph{touching a specific variable} is $O\bigl((e_{\mathrm{sup}}(k-1))^\ell / (n\ell)\bigr) = O(1/n)$ for fixed~$\ell$.  Summing over $\ell = 2, \ldots, L$ for any fixed~$L$: the expected number of cascade branches pruned at any given variable is~$O(1/n)$.

\emph{Supercriticality is preserved.}  The effective offspring distribution at each variable is $S_v - R_v$, where $S_v \sim \mathrm{Poi}(e_{\mathrm{sup}})$ is the support count and $R_v \geq 0$ is the number of automatically repaired clauses.  Since $\mathbb{E}[R_v] = O(1/n)$, the effective offspring mean is $e_{\mathrm{sup}} - O(1/n) > 1$ for~$n$ sufficiently large.  The branching process remains supercritical, so $q_{\mathrm{graph}} < 1$ and the cascade diverges with positive probability.

Therefore no finite-length path in $G_{\mathrm{SAT}}$ can flip $x_i$ while maintaining satisfiability.
\end{proof}

\begin{remark}[Cycle robustness]\label{rem:cycle-robust}
The proof of~(b) establishes cycle robustness via two independent mechanisms.  First, the first-passage property ensures the cascade explores a tree of distinct variables even on a graph with cycles.  Second, cycle encounters in support clauses produce \emph{deterministic} free repairs (the sign always cooperates), which can only help the adversary---but the help is negligible ($O(1/n)$ branches pruned per variable).  For cryptographic applications, KeyGen can additionally verify in polynomial time that the designated frozen coordinate~$i$ has a cycle-free cascade neighborhood (only $O(1)$ supported variables fail this check, while $\Theta(fn)$ are available), providing a belt-and-suspenders guarantee.  Even without this check, the stochastic domination argument shows $f_{\mathrm{graph}} = f_{\mathrm{tree}} - O(1/n)$, so Theorem~\ref{thm:fci} holds unconditionally on the full graph.
\end{remark}

\begin{remark}[Cascade example]
Consider $k = 3$ with $y^* = (1,1,1,1,1,1)$ and clauses $c_1 = (x_1 \vee \neg x_2 \vee \neg x_3)$, $c_2 = (x_2 \vee \neg x_4 \vee \neg x_5)$, $c_3 = (x_3 \vee \neg x_5 \vee \neg x_6)$.  Under $y^*$, variable $x_1$ is the unique satisfier of $c_1$ (the negated literals $\neg x_2, \neg x_3$ are false).  Flipping $x_1$ breaks $c_1$; repairing $c_1$ by flipping $x_2$ (making $\neg x_2$ true) breaks $c_2$ (where $x_2$ was the unique satisfier); repairing $c_2$ by flipping $x_4$ breaks further clauses.  Each repair triggers new breaks---the cascade propagates through the factor graph.
\end{remark}

\begin{remark}\label{rem:params}
For $k = 7$ and $\alpha/\alpha_s = 0.78$: $e_{\mathrm{sup}} = 3.81$, the extinction probability is $q = 0.024$, and the frozen fraction is $f \approx 0.976$.  Nearly every variable is frozen.
\end{remark}


\subsection{Step 4: Hamming ball containment}

\begin{lemma}\label{lem:hamming}
Every solution $s \in \mathrm{Comp}(y^*)$ satisfies $\mathrm{Ham}(s, y^*) \leq (1 - f + \varepsilon)n$ for any $\varepsilon > 0$, w.h.p.
\end{lemma}

\begin{proof}
By Lemma~\ref{lem:frozen}(b), $s$ agrees with $y^*$ on all $fn - o(n)$ frozen coordinates.  Therefore $s$ and $y^*$ differ on at most $(1 - f + \varepsilon)n$ variables.
\end{proof}


\subsection{Step 5: Component disjointness}

\begin{proof}[Proof of Theorem~\ref{thm:fci}]
Item~(i) is Lemma~\ref{lem:support}.  Item~(ii) is Lemma~\ref{lem:frozen}.  Item~(iii) is Lemma~\ref{lem:hamming}.

For item~(iv): suppose $\mathrm{Ham}(y_1, y_2) > 2(1-f)n$.  By Lemma~\ref{lem:hamming}, $\mathrm{Comp}(y_1) \subset B(y_1, (1-f+\varepsilon)n)$ and $\mathrm{Comp}(y_2) \subset B(y_2, (1-f+\varepsilon)n)$.  By the triangle inequality, $B(y_1, r) \cap B(y_2, r) = \emptyset$ when $2r < \mathrm{Ham}(y_1, y_2)$.  Taking $r = (1-f+\varepsilon)n$ with $\varepsilon$ small, the condition $2(1-f+\varepsilon)n < \mathrm{Ham}(y_1, y_2)$ holds by hypothesis.
\end{proof}

\begin{remark}[Generality]\label{rem:general}
The proof uses only three structural properties: (i)~a computable support count exceeding~1, (ii)~sparse cycle density in the cascade subgraph (automatic for random hypergraphs at sub-satisfiability density), (iii)~disagreeing frozen cores between clusters.  Any CSP with these properties---including planted $k$-XORSAT, planted graph coloring, and planted $k$-NAE-SAT---satisfies an analogous isolation theorem.
\end{remark}


%% ===================================================================
\section{Corollary~\ref{cor:qi}: Quantum Isolation}
%% ===================================================================

\begin{proof}[Proof of Corollary~\ref{cor:qi}]
By Theorem~\ref{thm:fci}, $G_{\mathrm{SAT}}$ decomposes into connected components $G_1, \ldots, G_m$ respecting cluster boundaries.  Define $\mathcal{H}_i = \mathrm{span}\{|s\rangle : s \in G_i\}$.  The adjacency matrix $A$ is block-diagonal: $A = \bigoplus_i A_i$.  For any analytic $g$, $g(A) = \bigoplus_i g(A_i)$.  This includes $e^{-iAt}$, any polynomial $p(A)$, and $e^{-A^2 t}$.  If $|\psi(0)\rangle \in \mathcal{H}_i$, then $g(A)|\psi(0)\rangle \in \mathcal{H}_i$.
\end{proof}

\begin{remark}
This corollary is mathematically elementary; the contribution is the observation that frozen-variable structure in planted CSPs implies quantum isolation, a connection that appears to be new.  It rules out any unitary generated by the solution-subgraph structure, not only time evolution.  Algorithms operating on the full $2^n$-dimensional Hilbert space are addressed in \S\ref{sec:quantum}.
\end{remark}


%% ===================================================================
\section{Commitment Scheme}
%% ===================================================================

\subsection{Construction}

The scheme instantiates a standard sigma protocol~\cite{GMW91} with planted $k$-SAT.

\medskip\noindent\textbf{KeyGen.}  Generate a planted instance $\Phi$ with solution $y^*$.  Identify a frozen coordinate $i$ (verifiable by the key generator, who knows $y^*$, by checking that $x_i$ has $\geq 1$ support clause).  Public key: $(\Phi, i)$.  Secret key: $y^*$.

\medskip\noindent\textbf{Commit($b$).}  Produce a non-interactive zero-knowledge proof that ``I know $y$ satisfying $\Phi$ with $y_i = b$,'' using the following sigma protocol repeated $\lambda$ times with Fiat--Shamir~\cite{FS86}:
\begin{enumerate}
\item Pick random permutation $\pi$ on $[n]$ and random sign flip $\sigma \in \{-1,+1\}^n$.
\item Compute $\Phi' = (\pi, \sigma)(\Phi)$ and $y' = (\pi, \sigma)(y^*)$.
\item Send commitment $c = H(\Phi' \| y')$.
\item On challenge $b_c \in \{0,1\}$: if $b_c = 0$, reveal $(\pi,\sigma)$; if $b_c = 1$, reveal $y'$.
\end{enumerate}

\medskip\noindent\textbf{Open($b$).}  Reveal $b$ and the proof transcript.

\medskip\noindent\textbf{Verify.}  Check the proof transcript.


\subsection{Security proofs}

\begin{proposition}[Completeness]\label{prop:complete}
An honest prover succeeds with probability~$1$.
\end{proposition}

\begin{proof}
On $b_c = 0$: the prover reveals $(\pi,\sigma)$; the verifier recomputes $(\pi,\sigma)(\Phi)$ and confirms it matches the committed $\Phi'$.  This always holds by construction.  On $b_c = 1$: the prover reveals $y' = (\pi,\sigma)(y^*)$.  Since $y^*$ satisfies $\Phi$ and $(\pi,\sigma)$ preserves satisfiability, $y'$ satisfies $\Phi'$.  The verifier confirms this.
\end{proof}

\begin{proposition}[Soundness]\label{prop:sound}
A prover who does not know any solution to $\Phi$ succeeds with probability at most $2^{-\lambda}$.
\end{proposition}

\begin{proof}
If the prover can answer both challenges for a single commitment, then from $b_c = 0$ we extract $(\pi,\sigma)$ and from $b_c = 1$ we extract $y'$.  Computing $(\pi,\sigma)^{-1}(y')$ yields a solution to $\Phi$---contradicting the assumption.  Therefore the prover can answer at most one challenge per round.  Over $\lambda$ independent rounds, the success probability is at most $2^{-\lambda}$.
\end{proof}

\begin{proposition}[Zero-knowledge]\label{prop:zk}
There exists a polynomial-time simulator producing transcripts computationally indistinguishable from real transcripts (under collision-resistance of $H$).
\end{proposition}

\begin{proof}
The simulator guesses the challenge $b'_c$ uniformly.  If $b'_c = 0$: pick random $(\pi,\sigma)$, set $\Phi' = (\pi,\sigma)(\Phi)$, commit with a random dummy $y_{\mathrm{fake}}$; if the actual challenge is $b_c = 0$, reveal $(\pi,\sigma)$ (correct).  If $b'_c = 1$: generate a random planted instance $\Phi'$ with a fresh planted solution $y'$, commit $H(\Phi' \| y')$; if $b_c = 1$, reveal $y'$ (correct).  On mismatch ($b_c \neq b'_c$), abort and restart.  Expected restarts per round: $2$.

The simulated transcripts are identically distributed to real transcripts.  On $b_c = 0$: both real and simulated reveal a uniform $(\pi,\sigma)$.  On $b_c = 1$: the real transcript reveals $(\pi,\sigma)(y^*)$, a random solution to a random instance $(\pi,\sigma)(\Phi)$; the simulator reveals a random solution to a random planted instance.  These are identically distributed because the planted model is invariant under $(\pi,\sigma)$.
\end{proof}


\subsection{Binding via frozen-core isolation}\label{sec:binding-empirical}

\noindent\textbf{Binding.}  Suppose an adversary opens to both $b = 0$ and $b = 1$.  Then the adversary possesses solutions $y_0, y_1$ to $\Phi$ with $(y_0)_i = 0$ and $(y_1)_i = 1$.  Since $i$ is frozen, all solutions in $\mathrm{Comp}(y^*)$ satisfy $y_i = y^*_i$ (Definition~\ref{def:frozen}).  At least one of $y_0, y_1$ lies outside $\mathrm{Comp}(y^*)$.  Finding such a solution requires solving planted $k$-SAT independently.  Reduction: break binding $\Rightarrow$ solve planted $k$-SAT.

\begin{remark}
Binding and hiding require only that $\Phi$ is computationally hard to solve, not that $\Phi$ is indistinguishable from random.  For applications requiring indistinguishability (CCA security), parameters should be chosen in the quiet planting window~\cite{ACO08}; for $k \geq 15$ this window overlaps with the frozen regime.
\end{remark}

\begin{remark}[Empirical binding verification]\label{rem:binding-empirical}
We verified the binding structure computationally at $n = 20$--$75$.  For each supported variable $x_i$ (i.e., $x_i$ has $\geq 1$ support clause), we used Glucose to search for a satisfying assignment with $x_i$ forced to $\neg y^*_i$.  The solver succeeds at all sizes tested, but the alternative solutions lie in a \emph{different cluster}: the minimum Hamming distance from $y^*$ satisfies $d_{\min}/n \geq 0.34$ at every size tested, with average $d/n \approx 0.45$ (Table~\ref{tab:separation}).  This confirms the binding mechanism: flipping a supported variable within $\mathrm{Comp}(y^*)$ is blocked by the cascade structure; the only route to a solution with $x_i \neq y^*_i$ is to solve planted $k$-SAT independently and land in a distant cluster.  The cost of this independent solve is the classical attack cost $2^{0.234n}$ established in \S\ref{sec:concrete}.
\end{remark}

\begin{table}[ht]
\centering
\caption{Cluster separation distance under forced variable flip.  For each $n$, 5 supported variables are flipped; the solver's alternative solution is measured against $y^*$.  Separation is $\Omega(n)$: $d_{\min}/n \geq 0.34$ at all sizes.}\label{tab:separation}
\begin{tabular}{@{}cccccc@{}}
\toprule
$n$ & Tests & $d_{\min}$ & $d_{\mathrm{avg}}$ & $d_{\min}/n$ & $d_{\mathrm{avg}}/n$ \\
\midrule
20 & 5 & 8  & 9.6  & 0.400 & 0.480 \\
30 & 5 & 16 & 17.0 & 0.533 & 0.567 \\
40 & 5 & 18 & 20.8 & 0.450 & 0.520 \\
50 & 5 & 17 & 21.4 & 0.340 & 0.428 \\
60 & 5 & 26 & 27.4 & 0.433 & 0.457 \\
75 & 3 & 28 & 33.3 & 0.373 & 0.444 \\
\bottomrule
\end{tabular}
\end{table}


\subsection{Concrete parameter selection}\label{sec:params}

We recommend $k = 7$ (minimum $k$ for strong freezing) at operating density $\alpha/\alpha_s = 0.78$ (deeply supercritical, $e_{\mathrm{sup}} = 3.81$, frozen fraction $f \approx 0.976$).  The variable count $n$ is sized to resist both classical cryptanalysis (\S\ref{sec:concrete}) and Grover search (\S\ref{sec:grover}).  Empirical benchmarking of Glucose~4 and MiniSat on our exact distribution yields a classical attack cost scaling as $2^{0.23n}$ operations (\S\ref{sec:concrete}); this is the binding constraint, requiring $n \geq \lambda / 0.234$.

\begin{table}[ht]
\centering
\caption{Concrete parameters at NIST security levels.  Public key uses seed-based compression (\S\ref{sec:seed}).  Parameter $n$ is sized by the classical CDCL attack cost $2^{0.234n}$ (\S\ref{sec:concrete}), which dominates the Grover bound $2^{0.488n}$ at these sizes.}\label{tab:params}
\begin{tabular}{@{}lcccccccc@{}}
\toprule
Level & $\lambda$ & $n$ & $k$ & $\alpha/\alpha_s$ & $m$ (clauses) & PK size & Proof size \\
\midrule
NIST-1 (128-bit) & 128 & 547 & 7 & 0.78 & 37{,}854 & 34\,B & 12.5\,KB \\
NIST-3 (192-bit) & 192 & 820 & 7 & 0.78 & 56{,}747 & 34\,B & 25.1\,KB \\
NIST-5 (256-bit) & 256 & 1{,}094 & 7 & 0.78 & 75{,}708 & 34\,B & 42.0\,KB \\
\bottomrule
\end{tabular}
\end{table}

\noindent\textbf{Public key size.}  With seed-based key compression (\S\ref{sec:seed}), the public key is a 32-byte seed plus a 2-byte frozen coordinate index, totaling \textbf{34~bytes} at all security levels.  Without compression, each clause stores $k$ variable indices ($\lceil \log_2 n \rceil$ bits each) plus $k$ sign bits, yielding 356--776\,KB; see Table~\ref{tab:compress} for the compression ratios.

\noindent\textbf{Commitment size.}  One SHA-256 hash $= 32$ bytes.

\noindent\textbf{Proof size.}  The sigma protocol runs $\lambda$ rounds.  Each round: one 32-byte commitment hash plus a response (either a 32-byte seed encoding $(\pi,\sigma)$ or an $n$-bit assignment $y'$).  Total: approximately $\lambda \cdot (32 + n/8)$ bytes.

\noindent\textbf{Comparison.}  The 34-byte public key is smaller than CRYSTALS-Kyber ($1{,}568$~bytes at NIST-1).  Proof/signature sizes (12.5--42\,KB) are comparable to SPHINCS+ (7.9--49.9\,KB).  The tradeoff is verification time: reconstructing $\Phi$ from the seed costs 0.4--1.4\,s (Table~\ref{tab:compress}), compared to microsecond verification for lattice-based schemes.  For commitment and signature applications where signing/verification latency is not critical, this overhead is acceptable.


\subsection{Seed-based key compression}\label{sec:seed}

The na\"{i}ve public key---the full clause list---is large (160--342\,KB).  We eliminate this overhead entirely by deriving the instance deterministically from a short seed.

\begin{enumerate}
\item Sample a uniform seed $s \leftarrow \{0,1\}^{256}$.
\item Derive the planted solution: $y^* = \mathrm{SHAKE256}(s \,\|\, \texttt{"sol"}, \; n \text{ bits})$.
\item Derive clause randomness: $r = \mathrm{SHAKE256}(s \,\|\, \texttt{"cls"}, \; \text{sufficient length})$.
\item Generate $\Phi$ from $r$ using the standard planted model (variable selection by Fisher--Yates shuffle with bytes from $r$; literal signs sampled from $r$ conditioned on $y^*$ satisfying each clause).
\item Identify a frozen coordinate $i$: the first variable with $\geq 1$ support clause.
\item Public key: $(s, i)$.  Secret key: $y^*$.
\end{enumerate}

The verifier reconstructs $\Phi$ from $s$ by executing the identical deterministic procedure.  Both parties obtain the same instance; the clause list is never transmitted.

\begin{proposition}[Security of seed-based compression]\label{prop:seed}
If $\mathrm{SHAKE256}$ is a secure pseudorandom function, then the seed-based scheme is computationally indistinguishable from the scheme with true randomness.
\end{proposition}

\begin{proof}
Suppose a distinguisher $\mathcal{D}$ can distinguish a planted instance $\Phi_{\mathrm{seed}}$ generated from $\mathrm{SHAKE256}(s)$ from an instance $\Phi_{\mathrm{rand}}$ generated with true randomness, with advantage $\varepsilon$.  Then $\mathcal{D}$ distinguishes the output of $\mathrm{SHAKE256}$ from a uniform random string of the same length with the same advantage $\varepsilon$, contradicting pseudorandomness.  All security properties of the commitment scheme (binding, hiding, soundness, zero-knowledge) transfer from the true-randomness setting to the seeded setting with at most negligible loss.
\end{proof}

This is exactly the technique used by CRYSTALS-Kyber, whose public matrix $\mathbf{A}$ is derived from a 32-byte seed $\rho$ via SHAKE128~\cite{DSS22}.

\begin{table}[ht]
\centering
\caption{Seed-based key compression: empirical validation.}\label{tab:compress}
\begin{tabular}{@{}lccccc@{}}
\toprule
Level & $n$ & Uncompressed PK & Compressed PK & Ratio & Reconstruction \\
\midrule
NIST-1 & 547 & 364{,}345\,B & 34\,B & $10{,}716\times$ & 379\,ms \\
NIST-3 & 820 & 546{,}190\,B & 34\,B & $16{,}064\times$ & 840\,ms \\
NIST-5 & 1{,}094 & 794{,}934\,B & 34\,B & $23{,}380\times$ & 1{,}424\,ms \\
\bottomrule
\end{tabular}
\end{table}

At every security level tested, deterministic reconstruction from the 34-byte public key produces a bit-identical instance, with the empirical average support per variable ($3.61$--$3.86$) matching the theoretical prediction $e_{\mathrm{sup}} = 3.81$.


\subsection{Digital signature scheme}\label{sec:sig}

The sigma protocol of \S5.1 yields a digital signature scheme via the standard Fiat--Shamir transform~\cite{FS86}, binding the message into the challenge derivation.

\medskip\noindent\textbf{Sign($\mathrm{msg}$, $y^*$).}  For $j = 1, \ldots, \lambda$:
\begin{enumerate}
\item Pick random $(\pi_j, \sigma_j)$.  Compute $\Phi'_j = (\pi_j, \sigma_j)(\Phi)$ and $y'_j = (\pi_j, \sigma_j)(y^*)$.
\item Compute commitment $a_j = H(\Phi'_j \| y'_j)$.
\end{enumerate}
Derive challenge bits $(b_1, \ldots, b_\lambda) = H(a_1 \| \cdots \| a_\lambda \| \mathrm{msg})$.  For each $j$: if $b_j = 0$, set $z_j = (\pi_j, \sigma_j)$; if $b_j = 1$, set $z_j = y'_j$.  The signature is $\sigma_{\mathrm{sig}} = (a_1, \ldots, a_\lambda, z_1, \ldots, z_\lambda)$.

\medskip\noindent\textbf{Verify($\mathrm{msg}$, $\sigma_{\mathrm{sig}}$, $(s, i)$).}  Reconstruct $\Phi$ from seed $s$.  Recompute $(b_1, \ldots, b_\lambda) = H(a_1 \| \cdots \| a_\lambda \| \mathrm{msg})$.  For each $j$: if $b_j = 0$, check that $(\pi_j, \sigma_j)(\Phi)$ is consistent with $a_j$; if $b_j = 1$, check that $y'_j$ satisfies $\Phi'_j$ and is consistent with $a_j$.

\begin{proposition}[Unforgeability]\label{prop:sig}
The signature scheme is existentially unforgeable under chosen-message attack (EUF-CMA) in the random oracle model, assuming planted $k$-SAT is hard.
\end{proposition}

\begin{proof}
By the Fiat--Shamir paradigm applied to a proof of knowledge~\cite{FS86,GMW91}: any forger producing a valid signature on a new message can be rewound (by reprogramming the random oracle) to extract two accepting transcripts for the same commitments with different challenges.  From these, as in Proposition~\ref{prop:sound}, one extracts a solution to $\Phi$.  Therefore forging a signature reduces to solving planted $k$-SAT.
\end{proof}

\noindent\textbf{Signature size.}  Each round contributes 32~bytes (commitment) plus a response: either a 32-byte seed encoding $(\pi_j, \sigma_j)$ or an $n$-bit string $y'_j$.  Total: approximately $\lambda \cdot (32 + n/8)$ bytes---identical to the proof size in Table~\ref{tab:params}.


%% ===================================================================
\section{Quantum Attack Analysis}\label{sec:quantum}
%% ===================================================================

We prove exponential lower bounds against five quantum algorithmic families.  The sixth (Grover) is treated in detail in~\S\ref{sec:grover}.


\subsection{Quantum walks}
Corollary~\ref{cor:qi}: zero amplitude transfer between clusters for all $t \geq 0$.  Exact, not approximate.


\subsection{QAOA at all depths}
At constant depth $p = O(1)$, the light-cone has radius $O(k^p)$.  Clusters disagree on $\Omega(n)$ variables, so overlap is exponentially small.  At superconstant depth, the mixing operator restricted to the solution subspace is a polynomial in $A$; Corollary~\ref{cor:qi} applies.


\subsection{Quantum annealing}

\begin{proposition}\label{prop:anneal}
The minimum spectral gap of the adiabatic interpolation $H(s) = (1-s)H_0 + sH_P$ satisfies $\Delta_{\min} \leq e^{-\Omega(n)}$.
\end{proposition}

\begin{proof}
At $s = 1$, satisfying assignments in distinct components are degenerate ground states.  The transverse field $H_0$ lifts this degeneracy via tunneling.  To tunnel from $\mathrm{Comp}(y_1)$ to $\mathrm{Comp}(y_2)$, at least $\Omega(fn)$ frozen variables must be simultaneously flipped through a region where each contributes energy cost $\geq 1$ (one support clause violated per frozen variable).  By the Anderson localization bound~\cite{AKR10}, the tunneling amplitude is $e^{-\Omega(fn)}$ and $\Delta_{\min} \leq e^{-\Omega(n)}$.  The adiabatic runtime is $T \geq e^{\Omega(n)}$.
\end{proof}

\begin{remark}
The energy cost per frozen variable ($\geq 1$) exceeds any constant transverse-field strength $\Gamma$, confirming that the AKR localization conditions~\cite{AKR10} are satisfied at our parameters.
\end{remark}


\subsection{Montanaro DPLL walk}
Montanaro~\cite{Mon18} achieves quadratic speedup over classical backtracking.  Classical DPLL on planted $k$-SAT at our densities requires $2^{\Omega(n)}$ time~\cite{ACO08}.  Quantum: $2^{\Omega(n/2)}$---still exponential.


\subsection{Low-degree polynomials}
Bresler--Huang~\cite{BH22}: degree-$D$ polynomial algorithms cannot distinguish planted from random for $D = o(n)$.  This subsumes SOS relaxations, statistical queries, and constant-depth quantum circuits.


\subsection{Grover's algorithm}\label{sec:grover}

Grover's algorithm treats the solution space as an unstructured database, achieving quadratic speedup: $O(2^{n/2}/\sqrt{|S|})$ queries to an oracle that checks satisfiability.  For our commitment scheme, the relevant attack is: given $\Phi$ and frozen coordinate $i$, find any solution $y$ with $y_i \neq y^*_i$.  This requires finding a solution in a \emph{different} cluster from~$y^*$.

\paragraph{Effective search space.}
By Theorem~\ref{thm:fci}(iv), solutions outside $\mathrm{Comp}(y^*)$ lie at Hamming distance $> 2(1-f)n$ from $y^*$.  For $k = 7$, $\alpha/\alpha_s = 0.78$: $f \approx 0.976$, so the adversary must find a solution differing from $y^*$ in at least $95.2\%$ of variables.  Each such solution resides in a cluster whose discovery constitutes an independent planted $k$-SAT instance.  Grover provides no structural advantage beyond brute-force search over this effective space.

\paragraph{Query complexity lower bound.}
Grover's algorithm is optimal for unstructured search~\cite{BBBV97,Zal99}.  Any quantum algorithm making $T$ queries to a satisfiability oracle succeeds with probability at most $O(T^2 |S| / 2^n)$.  For our parameters, $|S|/2^n$ is exponentially small.  Since the adversary must find a solution outside $\mathrm{Comp}(y^*)$, and each cluster is confined to a Hamming ball of radius $(1-f)n$, the effective Grover complexity is $\Omega(2^{fn/2}) = \Omega(2^{0.488n})$.

\paragraph{Parameter sizing.}
For $\lambda$-bit security against Grover: set $n \geq 2\lambda/f$.  However, classical CDCL attacks scale as $2^{0.234n}$ (\S\ref{sec:concrete}), which is the binding constraint at all NIST levels.  The recommended parameters (Table~\ref{tab:params}) are sized by the classical bound $n \geq \lambda / 0.234$, which automatically satisfies the weaker Grover requirement $n \geq 2\lambda / f \approx 2.05\lambda$.  Compare: AES-128 requires doubling key length to 256~bits for Grover resistance; our classical bottleneck absorbs the quantum threat with margin.

\paragraph{Oracle cost.}
Each Grover oracle call requires evaluating $\Phi$ on a candidate assignment: checking $m = \alpha n$ clauses of width $k$.  This costs $O(k \alpha n) = O(n)$ classical gates per oracle call, or $O(n)$ Toffoli gates in a reversible circuit.  The total quantum circuit depth is $O(n \cdot 2^{fn/2})$, which is exponential.  The constant factors from reversible arithmetic are comparable to those in Grover-on-AES analyses~\cite{GLRS16}.

\paragraph{Standard treatment.}
Every deployed and proposed post-quantum scheme handles Grover identically: parameter doubling.  CRYSTALS-Kyber, CRYSTALS-Dilithium, SPHINCS+, McEliece---all set parameters to absorb the $\sqrt{\cdot}$ speedup.  No scheme has a structural Grover barrier because no such barrier exists for any NP search problem~\cite{BBBV97}.  Our treatment is therefore exactly on par with the NIST-standardized schemes.  The critical distinction is that the \emph{other} five quantum attack families (Table~\ref{tab:quantum}) \emph{do} admit structural barriers, which we prove.  Grover is the only residual threat precisely because it ignores all problem structure.


\subsection{Summary}

\begin{table}[ht]
\centering
\caption{Quantum attack analysis.  Five of six major quantum algorithmic families face proved structural barriers.  Grover's algorithm, the sole remaining threat, is mitigated by standard parameter doubling (\S\ref{sec:grover}), identically to how all NIST-standardized PQC schemes handle unstructured quantum search.}\label{tab:quantum}
\begin{tabular}{@{}llll@{}}
\toprule
Attack & Model & Barrier & Status \\
\midrule
Quantum walks     & $G_{\mathrm{SAT}}$ & Cor.~\ref{cor:qi} & Proved \\
QAOA (all depths) & $G_{\mathrm{SAT}}$ & Cor.~\ref{cor:qi} + light cone & Proved \\
Quantum annealing & $2^n$   & $\Delta \leq e^{-\Omega(n)}$ & Proved \\
Montanaro DPLL    & DPLL tree & $2^{\Omega(n/2)}$ & Proved \\
Low-degree poly   & Statistical & \cite{BH22} & Proved \\
Grover            & $2^n$   & $\Omega(2^{0.488n})$ & Mitigated \\
Arbitrary BQP     & $2^n$   & Assumption & Same as all PQ \\
\bottomrule
\end{tabular}
\end{table}

We emphasize that no existing post-quantum scheme---including CRYSTALS-Kyber, CRYSTALS-Dilithium, and SPHINCS+---provides a structural Grover barrier.  Our scheme matches this standard while providing proved barriers against five additional quantum attack families that lattice-based and hash-based schemes do not explicitly address.


%% ===================================================================
\section{Empirical Validation}
%% ===================================================================

\subsection{Multi-seed validation ($n = 16$)}

We generated planted instances with two solutions $y_1, y_2$ at Hamming distance $n/2$ across 100 random seeds at $n = 16$, $k = 7$, $\alpha/\alpha_s = 0.78$.  For each seed we enumerated all satisfying assignments, clustered by proximity, built the Hamming-1 solution subgraph, and checked for shared connected components between clusters.

\begin{table}[ht]
\centering
\caption{Multi-seed validation: $n = 16$, $k = 7$, $\alpha/\alpha_s = 0.78$, 100 seeds.}\label{tab:multiseed}
\begin{tabular}{@{}ll@{}}
\toprule
Metric & Value \\
\midrule
Valid instances (both clusters non-empty) & 100/100 \\
Min inter-cluster distance (every seed) & 16 ($= n$) \\
Seeds with shared $C_1$/$C_2$ components & 0/100 \\
Complete cluster isolation & 100/100 \\
\bottomrule
\end{tabular}
\end{table}

At every seed tested, the minimum inter-cluster Hamming distance equals $n$ and no connected component contains solutions from both clusters.


\subsection{Scaling ($n = 14$--$22$, single seed)}

\begin{table}[ht]
\centering
\caption{Barrier scaling at $k = 7$, $\alpha/\alpha_s = 0.78$, seed 42.}\label{tab:scaling}
\begin{tabular}{@{}ccccccl@{}}
\toprule
$n$ & $|S|$ & $|C_1|$ & $|C_2|$ & MinInter & Comp & Peak $|C_2|_2$ \\
\midrule
14 & 12 & 6 & 6 & 14 & 11 & 0 \\
16 & 13 & 7 & 6 & 16 & 11 & 0 \\
18 & 20 & 12 & 8 & 18 & 19 & 0 \\
20 & 27 & 11 & 16 & 20 & 25 & 0 \\
22 & 39 & 24 & 15 & 22 & 33 & 0 \\
\bottomrule
\end{tabular}
\end{table}


\subsection{Phase transition in $k$}

\begin{table}[ht]
\centering
\caption{Phase transition in $k$: $n = 20$, $\alpha/\alpha_s = 0.65$, seed 42.  Lower density than Table~\ref{tab:scaling} to show that barrier onset depends primarily on $k$.  The $k = 9$ entry has MinInter $= 2$ yet zero leakage, confirming that the mechanism is component fragmentation.}\label{tab:kphase}
\begin{tabular}{@{}ccccl@{}}
\toprule
$k$ & $|S|$ & MinInter & Comp & Peak $|C_2|_2$ \\
\midrule
5 & 82 & 3 & 25 & $4.7 \times 10^{-1}$ (Leaks) \\
7 & 27 & 20 & 25 & 0 (Blocked) \\
9 & 140 & 2 & 116 & 0 (Blocked) \\
\bottomrule
\end{tabular}
\end{table}


\subsection{Hop stress test}

\begin{table}[ht]
\centering
\caption{Hop sweep: $n = 22$, $k = 7$, $\alpha/\alpha_s = 0.78$.  Barrier holds through hop~4.}\label{tab:hop}
\begin{tabular}{@{}cccll@{}}
\toprule
Hop & Edges & Comp & Peak $|C_2|_2$ & \\
\midrule
1 & 15 & 30 & 0 & Blocked \\
2 & 27 & 29 & 0 & Blocked \\
4 & 27 & 29 & 0 & Blocked \\
5 & 66 & 12 & $3.4 \times 10^{-2}$ & Leaks \\
\bottomrule
\end{tabular}
\end{table}


\subsection{Concrete classical cryptanalysis}\label{sec:concrete}

To establish concrete bit-security, we benchmarked two state-of-the-art CDCL solvers---Glucose~4 and MiniSat~2.2---on planted instances at our exact operating parameters ($k = 7$, $\alpha/\alpha_s = 0.78$), averaging over 5~random seeds per size.  Instances were generated using the seed-based procedure of \S\ref{sec:seed}.

\begin{table}[ht]
\centering
\caption{CDCL solver performance on planted $k$-SAT at $k = 7$, $\alpha/\alpha_s = 0.78$.  Average over 5 seeds.  Timeout: 30\,s.}\label{tab:cryptanalysis}
\begin{tabular}{@{}cccccc@{}}
\toprule
$n$ & Clauses & Glucose time & Glucose conflicts & MiniSat time & MiniSat conflicts \\
\midrule
30  & 2{,}076  & 0.024\,s     & 2{,}559   & 0.007\,s     & 1{,}103 \\
40  & 2{,}768  & 0.070\,s     & 6{,}679   & 0.031\,s     & 4{,}245 \\
50  & 3{,}460  & 0.610\,s     & 32{,}678  & 0.381\,s     & 37{,}973 \\
60  & 4{,}152  & 2.628\,s     & 124{,}874 & 0.698\,s     & 61{,}192 \\
75  & 5{,}190  & $> 30$\,s    & 975{,}780 & $> 30$\,s    & 733{,}584 \\
\bottomrule
\end{tabular}
\end{table}

\paragraph{Exponential fit.}
Fitting $\log_2(\text{time}) = a + \beta n$ on the non-timeout data ($n = 30$--$60$) gives:
\begin{itemize}
\item Glucose: $\beta = 0.234$, $R^2 = 0.985$.
\item MiniSat: $\beta = 0.237$, $R^2 = 0.956$.
\end{itemize}
Fitting on conflict counts (a hardware-independent metric):
\begin{itemize}
\item Glucose conflicts: $\beta = 0.191$, $R^2 = 0.992$.
\item MiniSat conflicts: $\beta = 0.205$, $R^2 = 0.952$.
\end{itemize}
Taking the most conservative (attacker-favorable) fit, $\beta = 0.234$, the classical attack cost is $2^{0.234n}$ operations.

\paragraph{Extrapolated bit-security.}

\begin{table}[ht]
\centering
\caption{Extrapolated classical bit-security at recommended parameters.  Growth rate $2^{0.234n}$ from Glucose CDCL fit.}\label{tab:bitsec}
\begin{tabular}{@{}lcccc@{}}
\toprule
Level & $n$ & Classical cost & Grover cost & Security target \\
\midrule
NIST-1 & 547 & $2^{128}$ & $2^{267}$ & $2^{128}$ \\
NIST-3 & 820 & $2^{192}$ & $2^{400}$ & $2^{192}$ \\
NIST-5 & 1{,}094 & $2^{256}$ & $2^{534}$ & $2^{256}$ \\
\bottomrule
\end{tabular}
\end{table}

At the recommended parameters, the classical CDCL attack cost meets the NIST security target at each level, and the Grover (quantum brute-force) cost exceeds it by a wide margin.  The classical bottleneck dominates: parameters are sized by the $2^{0.234n}$ classical bound, not the $2^{0.488n}$ quantum bound.

\begin{remark}
The growth rate $\beta \approx 0.23$ is consistent with known hardness results for planted $k$-SAT in the frozen regime.  CDCL solvers are the strongest known general-purpose classical attack; Survey Propagation and Belief Propagation guided decimation are known to fail above the freezing threshold~\cite{ACO08}.  We are not aware of any classical algorithm achieving $\beta < 0.20$ on our distribution.
\end{remark}


%% ===================================================================
\section{Discussion}
%% ===================================================================

\subsection{Relationship to the Overlap Gap Property}

Theorem~\ref{thm:fci} is related to but distinct from Gamarnik's Overlap Gap Property (OGP)~\cite{Gam21}.  The OGP shows that solution spaces lack smooth interpolation paths between distant configurations.  Theorem~\ref{thm:fci} is strictly stronger in one sense: it proves topological disconnection, not merely absence of smooth paths.  However, the OGP applies to a broader class of problems.  Our result can be viewed as a rigorous strengthening of the OGP for planted $k$-SAT above the planted freezing threshold.


\subsection{A fourth hardness family}

Lattices, hash functions, codes, and now frozen-core constraint satisfaction.  The structural evidence for quantum resistance (Table~\ref{tab:quantum}) is, to our knowledge, more explicit in its coverage of specific quantum algorithmic families than what is available for any existing post-quantum scheme.  The recent proposal of quantum stabilizer decoding as a PQC hardness family~\cite{LPQR26} and the present work both suggest that the PQC assumption space is richer than the current NIST portfolio of lattices, hashes, and codes.  CSP cluster structure and quantum code decoding represent two independent directions for diversification.


\subsection{On the hardness assumption}

Our binding security reduces to the average-case hardness of planted $k$-SAT at specific parameters.  We acknowledge that a worst-case to average-case reduction---analogous to Regev's reduction from GapSVP to LWE---would strengthen the proposal, and we list this as an open problem.  We note, however, that this limitation is shared by multiple deployed and standardized PQC schemes:
\begin{itemize}
\item Hash-based signatures (SPHINCS+/FIPS~205) rest on the assumption that SHA-256 is collision-resistant.  No worst-case reduction to any NP-hard problem is known.
\item Code-based encryption (Classic McEliece) rests on the average-case hardness of decoding random linear codes.
\item Lattice-based schemes, while enjoying worst-case reductions in principle, derive their concrete parameter security from heuristic cost models of the BKZ algorithm at specific block sizes---an average-case extrapolation comparable to our CDCL benchmarks.
\end{itemize}
Our concrete security analysis (\S\ref{sec:concrete}) follows the same methodology used in the lattice PQC literature: benchmark the best known attack at small parameters, fit an exponential, and extrapolate.  The $R^2 = 0.985$ fit quality is comparable to or better than typical BKZ cost model fits.


\subsection{Generality}

By Remark~\ref{rem:general}, the isolation theorem applies to any CSP with frozen-variable structure.  If planted $k$-SAT is broken by a future SAT-specific algorithm, the framework survives with a different frozen-core CSP.


\subsection{Information-geometric connection}

The topological fragmentation proved here corresponds to a divergence of the Fisher--Rao metric on the Bernoulli manifold at the cluster boundary.  This connection will be developed in forthcoming work.


\subsection{Open problems}

\begin{enumerate}[label=(\roman*)]
\item Prove a full BQP lower bound for planted $k$-SAT.  No existing post-quantum scheme has such a proof; achieving one here would be a first.
\item Establish formal relationships (reductions or separations) between the planted $k$-SAT hardness assumption and existing PQC assumptions (LWE, SIS, syndrome decoding).
\item Extend the constructions to a full key encapsulation mechanism (KEM).  This requires a lossy mode for planted $k$-SAT instances, which connects to the quiet planting problem; for $k \geq 15$ this appears feasible.
\end{enumerate}


%% ===================================================================
%% REFERENCES
%% ===================================================================

\begin{thebibliography}{99}

\bibitem{ACO08}
D.~Achlioptas, A.~Coja-Oghlan, Algorithmic barriers from phase transitions, \emph{FOCS} (2008).

\bibitem{ART09}
D.~Achlioptas, F.~Ricci-Tersenghi, Random formulas have frozen variables, \emph{SIAM J.\ Computing} 39(1), 260--280 (2009).

\bibitem{AKR10}
B.~Altshuler, H.~Krovi, J.~Roland, Anderson localization makes adiabatic quantum optimization fail, \emph{PNAS} 107(28), 12446--12450 (2010).

\bibitem{BBBV97}
C.\,H.~Bennett, E.~Bernstein, G.~Brassard, U.~Vazirani, Strengths and weaknesses of quantum computing, \emph{SIAM J.\ Computing} 26(5), 1510--1523 (1997).

\bibitem{BL17}
D.\,J.~Bernstein, T.~Lange, Post-quantum cryptography, \emph{Nature} 549, 188--194 (2017).

\bibitem{BH22}
G.~Bresler, B.~Huang, Algorithmic phase transition of random $k$-SAT for low degree polynomials, \emph{FOCS} (2022).

\bibitem{Cho25}
K.-B.~Cho et~al., Equilibrium SAT based PQC, arXiv:2512.02598 (2025).

\bibitem{CEMT09}
J.~Cook, O.~Etesami, R.~Miller, L.~Trevisan, Goldreich's one-way function candidate and myopic backtracking algorithms, \emph{TCC 2009}, LNCS 5444, 521--538.

\bibitem{CZ12}
A.~Coja-Oghlan, L.~Zdeborov\'{a}, The condensation transition in random hypergraph 2-coloring, \emph{SODA} (2012).

\bibitem{DM10}
A.~Dembo, A.~Montanari, Gibbs measures and phase transitions on sparse random graphs, \emph{Brazilian J.\ Probability} 24(2), 137--211 (2010).

\bibitem{DSS22}
J.~Ding, A.~Sly, N.~Sun, Proof of the satisfiability conjecture for large $k$, \emph{Annals of Mathematics} 196(1), 1--388 (2022).

\bibitem{FGG14}
E.~Farhi, J.~Goldstone, S.~Gutmann, A quantum approximate optimization algorithm, arXiv:1411.4028 (2014).

\bibitem{FPV18}
V.~Feldman, W.~Perkins, S.~Vempala, On the complexity of random satisfiability problems with planted solutions, \emph{SIAM J.\ Computing} 47(4), 1294--1338 (2018).

\bibitem{FS86}
A.~Fiat, A.~Shamir, How to prove yourself, \emph{Crypto} (1986).

\bibitem{Gam21}
D.~Gamarnik, The overlap gap property, \emph{PNAS} 118(41) (2021).

\bibitem{GLRS16}
M.~Grassl, B.~Langenberg, M.~Roetteler, R.~Steinwandt, Applying Grover's algorithm to AES: quantum resource estimates, \emph{PQCrypto 2016}, LNCS 9606, 29--43.

\bibitem{Gol00}
O.~Goldreich, Candidate one-way functions based on expander graphs, ECCC TR00-090 (2000).

\bibitem{GMW91}
O.~Goldreich, S.~Micali, A.~Wigderson, Proofs that yield nothing but their validity, \emph{JACM} 38(3), 691--729 (1991).

\bibitem{LPQR26}
J.\,Z.~Lu, A.~Poremba, Y.~Quek, A.~Ramkumar, Post-quantum cryptography from quantum stabilizer decoding, ePrint 2026/548 (2026).

\bibitem{Mol12}
M.~Molloy, The freezing threshold for $k$-colourings of a random graph, \emph{STOC} (2012).

\bibitem{MR95}
M.~Molloy, B.~Reed, A critical point for random graphs with a given degree sequence, \emph{Random Structures \& Algorithms} 6(2--3), 161--179 (1995).

\bibitem{Mon18}
A.~Montanaro, Quantum walk speedup of backtracking algorithms, \emph{Theory of Computing} 14(15), 1--24 (2018).

\bibitem{MMRE23}
H.~Munoz-Bauza, H.~Munoz-Bauza, G.~Mossi, E.\,G.~Rieffel, Generating hard Ising instances with planted solutions using post-quantum cryptographic protocols, arXiv:2308.09704 (2023).

\bibitem{Sch15}
S.\,E.~Schmittner, A SAT-based public key cryptography scheme, arXiv:1507.08094 / ePrint 2015/771 (2015).

\bibitem{Zal99}
C.~Zalka, Grover's quantum searching algorithm is optimal, \emph{Phys.\ Rev.\ A} 60, 2746 (1999).

\bibitem{ZK07}
L.~Zdeborov\'{a}, F.~Krzakala, Phase transitions in the coloring of random graphs, \emph{Phys.\ Rev.\ E} 76, 031131 (2007).

\end{thebibliography}

\end{document}
